\section*{Bab \ref{ch:b2:ism-star-formation} --- Medium Antarbintang dan Pembentukan Bintang}

\begin{solution}{exo:b2:ism-star-formation:1}
$n = \num{2.7e-24}/\num{1.67e-24} = 1.6$ atom hidrogen per sentimeter kubik. Dalam udara,
$n = P/k_{\mathrm B}T = 101325/(\num{1.381e-23}\times288) = \qty{2.55e25}{m^{-3}}$,
yaitu \num{2.5e19} molekul per sentimeter kubik: sekitar $1.6\times10^{19}$ kali
lebih banyak.
\end{solution}

\begin{solution}{exo:b2:ism-star-formation:2}
$M_J \propto n^{-1/2}$: membagi $n$ dengan 100 mengalikan $M_J$ dengan 10, sehingga
$M_J = 54.6$ massa Matahari pada $10^2$ partikel per sentimeter kubik. Awan itu memuat
$\num{24000}/54.6 \approx 440$ \omterm{def:b2:ism-star-formation:jeans}{massa Jeans}: awan itu tidak mungkin stabil secara keseluruhan, dan pasti
terfragmentasi.
\end{solution}

\begin{solution}{exo:b2:ism-star-formation:3}
$c_s = \sqrt{\num{1.381e-23}\times10/(2.33\times\num{1.661e-27})} =
\qty{189}{m.s^{-1}}$ pada \qty{10}{K}; $\sqrt{10}$ kali lebih besar, \qty{597}{m.s^{-1}},
pada \qty{100}{K}. Menyeberangi \qty{0.1}{pc}: $\num{3.086e15}/189 = \qty{1.63e13}{s}$,
yaitu \qty{0.52}{Ma}.
\end{solution}

\begin{solution}{exo:b2:ism-star-formation:4}
$\rho = 2.33\times\num{1.661e-27}\times n$ dengan $n$ dalam \unit{m^{-3}}. Untuk $10^5$
partikel per sentimeter kubik, $\rho = \qty{3.87e-16}{kg.m^{-3}}$ dan
$t_{\mathrm{ff}} = \sqrt{3\pi/(32G\rho)} = \qty{3.38e12}{s} = \qty{0.107}{Ma}$; untuk
$10^2$, $\rho$ \num{1000} kali lebih kecil dan $t_{\mathrm{ff}}$ $\sqrt{1000} = 31.6$
kali lebih lama, \qty{3.38}{Ma}. \omterm{def:b2:ism-star-formation:molecular-cloud}{Inti padat} runtuh 32 kali lebih cepat.
\end{solution}

\begin{solution}{exo:b2:ism-star-formation:5}
Dengan $k_{\mathrm B}T/(\mu m_u) = c_s^2$, $(5k_{\mathrm B}T/G\mu m_u)^{3/2} =
5^{3/2}c_s^3G^{-3/2}$, yang menghasilkan rumus itu. Karena $M_J \propto \rho^{-1/2}$,
untuk membuat $M_J$ menjadi separuh diperlukan kerapatan empat kali lebih tinggi. Karena $M_J \propto T^{3/2}$,
penghangatan dari \qty{10}{K} menjadi \qty{20}{K} mengalikannya dengan $2^{3/2} = 2.83$.
\end{solution}

\begin{solution}{exo:b2:ism-star-formation:6}
$t_{\mathrm{KH}} = GM^2/(RL) = \num{6.674e-11}\times(\num{1.989e30})^2/(2\times
\num{6.96e8}\times2\times\num{3.83e26}) = \qty{2.48e14}{s}$, yaitu
\qty{7.8}{Ma}. Waktu ini lebih panjang daripada umur sebagian besar cakram (separuhnya hilang setelah
sekitar \qty{3}{Ma}): bintang itu masih mengerut ketika cakramnya lenyap.
\end{solution}

\begin{solution}{exo:b2:ism-star-formation:7}
Pada \qty{10}{K}, $c_s^3/G = 189^3/\num{6.674e-11} = \qty{1.01e17}{kg.s^{-1}}$, yaitu
$\num{1.01e17}\times\num{3.156e7}/\num{1.989e30} = \num{1.6e-6}$ massa Matahari per
tahun, dan satu massa Matahari memerlukan \qty{0.62}{Ma}. Pada \qty{20}{K} lajunya
$2^{3/2} = 2.83$ kali lebih tinggi, \num{4.5e-6} massa Matahari per tahun, dan satu massa
Matahari memerlukan \qty{0.22}{Ma}.
\end{solution}

\begin{solution}{exo:b2:ism-star-formation:8}
$r = \sqrt{L/(16\pi\sigma)}/T^2$ dengan $\sqrt{\num{3.83e26}/(16\pi\times\num{5.67e-8})}
= \qty{1.159e16}{m.K^2}$: pada \qty{1300}{K}, $r = \qty{6.86e9}{m} = \qty{0.046}{au}$,
jauh di dalam orbit Merkurius. Pada \qty{5.2}{au}, $T = 278/\sqrt{5.2} = \qty{122}{K}$.
\end{solution}

\begin{solution}{exo:b2:ism-star-formation:9}
$v = \sqrt{GM/r} = \sqrt{\num{6.674e-11}\times1.3\times\num{1.989e30}/\num{1.496e13}}
= \qty{3.40}{km.s^{-1}}$, dan $P = 2\pi r/v = \qty{2.77e10}{s} = 877$ tahun
(secara setara, $\sqrt{100^3/1.3}$ tahun). Dalam sejuta tahun gas itu telah berputar
$10^6/877 \approx \num{1140}$ kali: cakram itu tua dibandingkan dengan rotasinya sendiri,
dan itulah sebabnya cincin-cincinnya berupa lingkaran yang mulus.
\end{solution}

\begin{solution}{exo:b2:ism-star-formation:10}
$M = \int_0^R 4\pi r^2\rho_0R/r\,\mathrm dr = 2\pi\rho_0R^3$, dan di dalam $r$,
$m(r) = 2\pi\rho_0Rr^2$. Dengan $\mathrm dm = 4\pi\rho_0Rr\,\mathrm dr$,
\[
  U = -\int_0^R \frac{G\,m(r)}{r}\,\mathrm dm
    = -8\pi^2G\rho_0^2R^2\int_0^R r^2\,\mathrm dr = -\frac{8\pi^2}{3}G\rho_0^2R^5
    = -\frac23\,\frac{GM^2}{R},
\]
dengan $\rho_0 = M/(2\pi R^3)$. Nilainya $\tfrac{2/3}{3/5} = 1.11$ kali nilai bola seragam:
awan yang terpusat terikat lebih erat dan memerlukan energi termal \qty{11}{\percent} lebih
banyak untuk ditopang; awan itu lebih mudah runtuh.
\end{solution}

\begin{solution}{exo:b2:ism-star-formation:11}
Periodenya $2\pi/\Omega = \qty{6.28e14}{s}$, sekitar \qty{20}{Ma}. Dengan
$R = \qty{3.086e15}{m}$, $j = \Omega R^2 = \qty{9.52e16}{m^2.s^{-1}}$, dan
$r_c = j^2/(GM) = \num{9.07e33}/\num{1.327e20} = \qty{6.83e13}{m} = \qty{457}{au}$.
Dalam rotasi benda tegar, $j = \Omega r^2$ tumbuh sebanding dengan kuadrat jarak dari sumbu:
gas dari bagian dalam inti memiliki momentum sudut lebih kecil, dan mendarat lebih dekat ke bintang.
\end{solution}

\begin{solution}{exo:b2:ism-star-formation:12}
$\bar t = 2.46$, $\bar f = 0.588$; $\sum(t - \bar t)^2 = 12.612$,
$\sum(t - \bar t)(f - \bar f) = -2.018$, sehingga $b = 2.018/12.612 = 0.160$ per juta
tahun dan $a = 0.588 + 0.160\times2.46 = 0.982$. Garis itu mencapai $0.5$ pada
$(0.982 - 0.5)/0.160 = \qty{3.0}{Ma}$ dan nol pada $0.982/0.160 = \qty{6.1}{Ma}$, seperti pada
gambar. Sebuah fraksi tidak dapat turun di bawah nol: gugus yang lebih tua, seperti NGC 1960 pada
\qty{30}{Ma}, mempertahankan fraksi kecil di dekat nol yang tidak dapat digambarkan oleh garis lurus.
\end{solution}

\begin{solution}{pb:b2:ism-star-formation:1}
\textbf{1.} $R = 0.08\times\num{3.086e16} = \qty{2.469e15}{m}$.

\textbf{2.} $N = \qty{1.5e26}{m^{-2}}$, sehingga $2.8\times\num{1.661e-27}\times\num{1.5e26} =
\qty{0.698}{kg.m^{-2}}$.

\textbf{3.} $M = 0.698\times\pi R^2 = 0.698\times\num{1.915e31} = \qty{1.34e31}{kg}$,
yaitu $6.7$ massa Matahari.

\textbf{4.} $V = \tfrac43\pi R^3 = \qty{6.30e46}{m^3}$ dan
$\rho = M/V = \qty{2.12e-16}{kg.m^{-3}}$.

\textbf{5.} $n = \rho/(2.33m_u) = \qty{5.48e10}{m^{-3}}$, yaitu \num{5.5e4} partikel
per sentimeter kubik.

\textbf{6.} Ya: kerapatannya terletak antara $10^4$ dan $10^5$, jari-jarinya antara 0.03 dan
\qty{0.2}{pc}; massanya, 6.7 massa Matahari, berada di ujung atas rentang inti.

\textbf{7.} $M_J = 5.46\times(10^4/\num{5.48e4})^{1/2} = 2.33$ massa Matahari.

\textbf{8.} $6.7/2.33 = 2.9$ \omterm{def:b2:ism-star-formation:jeans}{massa Jeans}: inti itu tidak stabil dan runtuh.

\textbf{9.} $R_J = 0.1316\times(10^4/\num{5.48e4})^{1/2} = \qty{0.056}{pc}$, lebih kecil daripada
jari-jarinya, \qty{0.08}{pc}: kesimpulannya sama.

\textbf{10.} $2^{3/2}\times2.33 = 6.6$ massa Matahari: sekitar massa inti itu. Pada
\qty{20}{K} inti itu berada di ambang; dinginlah yang membuatnya runtuh.

\textbf{11.} $t_{\mathrm{ff}} = \sqrt{3\pi/(32\times\num{6.674e-11}\times\num{2.12e-16})}
= \qty{4.56e12}{s}$, sekitar \num{145000} tahun.

\textbf{12.} $R/c_s = \num{2.469e15}/189 = \qty{1.31e13}{s}$, \num{414000} tahun: tiga
kali lebih lama daripada jatuh bebas. Gelombang tekanan tidak dapat menyeberangi inti sebelum inti itu
runtuh; tekanan tidak dapat bereaksi cukup cepat untuk menghentikan keruntuhan.

\textbf{13.} $c_s = \qty{189}{m.s^{-1}}$ (\cref{exo:b2:ism-star-formation:3}).

\textbf{14.} $c_s^3/G = \qty{1.01e17}{kg.s^{-1}}$, \num{1.6e-6} massa Matahari per tahun.

\textbf{15.} $1/\num{1.6e-6} = \num{624000}$ tahun.

\textbf{16.} Sepertiga dari 6.7 adalah $2.2$: dua bintang bermassa satu massa Matahari.

\textbf{17.} $t_{\mathrm{KH}} = GM_\odot^2/(3R_\odot\times4L_\odot) = \qty{8.25e13}{s}$,
\qty{2.6}{Ma}.

\textbf{18.} Lebih panjang, empat kali waktu akresi: ketika akresi berhenti, bintang itu
masih mengerut, bersinar tanpa fusi; ia adalah \omterm{def:b2:ism-star-formation:t-tauri}{bintang T Tauri}, belum berada di deret
utama.

\textbf{19.} $\Omega = \qty{0.4e3}{m.s^{-1}}/\qty{3.086e16}{m} = \qty{1.30e-14}{rad.s^{-1}}$.

\textbf{20.} $j = 1.30\times10^{-14}\times(\num{2.469e15})^2 = \qty{7.90e16}{m^2.s^{-1}}$.

\textbf{21.} $r_c = j^2/(GM_\odot) = \num{6.24e33}/\num{1.327e20} = \qty{4.70e13}{m}$, yaitu
\qty{314}{au}: lebih besar daripada cakram HL Tauri, seperti yang diharapkan untuk gas dari tepi
inti.

\textbf{22.} $T = (L/16\pi\sigma r^2)^{1/4} = \qty{278}{K}$ pada \qty{1}{au}, dan
$278/\sqrt{30} = \qty{51}{K}$ pada \qty{30}{au}.

\textbf{23.} $r = \num{1.159e16}/180^2 = \qty{3.58e11}{m}$, yaitu \qty{2.4}{au}, seperti dalam teks.

\textbf{24.} Bintang bermassa satu massa Matahari memerlukan
\textbf{\qty[text-series-to-math=true]{624}{} ribu tahun} untuk terhimpun, sekitar empat
kali \omterm{def:b2:ism-star-formation:free-fall}{waktu jatuh bebas} intinya (\num{145000} tahun): gas di bagian luar mencapai
\omterm{def:b2:ism-star-formation:protostar}{protobintang} hanya secara bertahap, dari dalam ke luar.
\end{solution}
